\section{Source-grounded definitions help, but text and scope matter}
\label{sec:repair}
\paragraph{Source content, not a prior-only gloss, drives the synthetic gain.}
On 761 synthetic questions whose drafter saw usage sentences but not a definition, DeepSeek V4 Flash accuracy is 0.699 without a note, 0.966 with the source definition, 0.949 with a model-drafted gloss \emph{based on that source}, and 0.708 with the reader's own domain-conditioned gloss (Figure~\ref{fig:silent}c). The source definition repairs 209 errors and harms 6 previously correct answers; the prior-only gloss repairs 86 and harms 79. A model can therefore write a useful paraphrase when it has the source, but asking it to supply a missing local rule from its own prior does not reproduce that benefit. In the 180 questions where a lexical-overlap screen singles out no option from any of the three notes, source-definition accuracy still rises from 0.617 to 0.911, versus 0.544 with the prior gloss. This check reduces one answer-cue concern without removing all construction effects (Appendix~\ref{app:lexical-cues}). For a semantic layer, the object to preserve is the source-bound rule and its scope, not merely a fluent gloss of the term.

The ability to use a supplied rule is separately testable. On 1,244 human-written Inverse Scaling redefinition items \citep{inverse-scaling}, GPT-5.6-terra and GPT-6-astra follow the stated meaning on 95.0\% and 99.8\%; without it both use the usual meaning on about 98\%. DeepSeek V4.1 Flash follows the stated meaning on 63.0\%. The contrast separates failure to \emph{receive} a local rule from failure to \emph{follow} one, and shows why the two should not be collapsed into a single intervention claim.

\paragraph{A definition's wording changes both repairs and harms.}
On Reddit glossary questions, a model-drafted source-based gloss repairs 171 errors but harms 58 correct answers; the community entry repairs 243 and harms 9. The corresponding DeFi counts are 42/20 and 70/3. In both corpora the question writer saw the source entry, so this is an informative text-version comparison, not an independent estimate of definition quality. It warns that supplying an inaccurate or poorly scoped note can make an answer worse. For the narrower target of \emph{which known bare errors the same note repairs}, $\Delta p$ ranks outcomes with AUROC 0.712--0.983 across eight heterogeneous settings (Appendix~\ref{app:statistics}). That diagnostic is tied to the exact text and the evaluated question. On 215 Reddit terms, a drafted-gloss score predicts sibling-question repairs at 0.811 AUROC, while a score computed with the community entry reaches 0.600 $[0.475,0.717]$; the latter entry already repairs 126 of 147 bare errors, leaving only 21 nonrepairs to rank (Appendix~\ref{app:new}). Neither result establishes that $\Delta p$ can choose definitions for unseen user requests.

\paragraph{A definition is not always sufficient.}
In CUAD, extracted contract definitions repair 163 errors but harm 83 correct answers, raising accuracy from 0.787 to 0.840 (Figure~\ref{fig:strata}b). Of 320 original errors, 157 remain. A definition may depend on schedules or other clauses, the question may be defective, or the reader may still apply it incorrectly. An operational system should therefore retrieve the linked evidence and allow an unresolved case to remain unresolved; simply attaching every matched term note is not a guarantee of correctness.

\paragraph{Which definitions to curate first.}
The cost of a definition library lies in writing, checking, and maintaining entries. In a 215-term Reddit development comparison at 25\% of total definition words, frequency-based selection has point estimates 3.5 and 4.5 accuracy points above frequency times mean $\Delta p$ for the two readers on policy-disagreement terms. Source-check uncertainty gives a difference interval of $[-9.21,+6.14]$ points, so neither policy is established as better. A small CUAD comparison also does not favour mean $\Delta p$ (Appendix~\ref{app:new}). The historical retrieval gains cannot establish advance selection because the evaluated question contributed to its own term score and the CUAD leave-one-question-out gain is only partially identified (Appendix~\ref{app:interventions}). For now, term frequency is a transparent default for deciding what to curate first; $R$ can separately order diagnostic questions for review. These are different decisions and should not be conflated.

